\section{Staticity, recovery of the original data, and sharpness}
\label{ri:section}

Throughout the rigidity argument the hypotheses of
Theorem~\ref{thm:rigidity} hold.  Set
\[
 A=4\pi r_h^2,\qquad r_h=m+\sqrt{m^2-Q^2}>Q,
 \qquad \kappa=\frac{1-Q^2/r_h^2}{2r_h}>0.
\]
We use the stationary metric $\mathbf G$, Maxwell form $\mathbf F$,
and Killing field $\xi=\partial_z$ constructed in
Proposition~\ref{va:stationary-data}.  In particular, their restriction
to the original section $z=0$ induces the original four tensors,
with the prescribed future normal and orientation.  The constructions
of Section~\ref{st:section} have not yet asserted staticity or removed
the possible residual matter stress.  Both conclusions will follow
from an integral on the maximal hypersurface.

\subsection{A global maximal graph and the staticity identity}
\label{ri:staticity-subsection}

The maximal-slice integral strategy below adapts the staticity method of
Sudarsky and Wald; compare \cite[Eq.~(40) and Theorem~2]{SudarskyWald1993}.
We derive the identity locally, including the duality rotation, global
magnetic-potential argument and aligned residual-matter term.  These
features and the present hypotheses are not supplied by the cited
electrovacuum staticity theorem.

\begin{lemma}[The maximal hypersurface as an orbit section]
\label{ri:orbit-graph}
Let $\Sigma'$ be the maximal hypersurface of
Proposition~\ref{st:maximal-slice}, and let $B_0$ be its bifurcation
boundary.  Projection along the Killing flow identifies
$\operatorname{int}\Sigma'$ diffeomorphically with
$\operatorname{int}\Omega$.  This identification takes the end to the
same end and identifies the boundary labels on $B_0$ with those on
$S$.  Compact cuts contained in the interiors of the two hypersurfaces,
with corresponding orbit labels, are homologous in the right wedge.
\end{lemma}

\begin{proof}
Each Killing orbit in the right wedge is a complete, nonzero causal
curve.  It is inextendible there: in the original product its spatial
label is fixed and its $z$ coordinate traverses all of $\mathbb R$.
The Cauchy property of $\operatorname{int}\Sigma'$ therefore gives
exactly one intersection with each orbit.  A nonzero causal vector
cannot be tangent to a spacelike hypersurface.  Consequently the
restriction of orbit projection is a local diffeomorphism, and its
bijectivity makes it a global diffeomorphism.  Flowing the section
thus gives a diffeomorphism
\[
 \mathbb R\times\operatorname{int}\Sigma'
       \longrightarrow \mathbb R\times\operatorname{int}\Omega.
\]
For a compact cut its two sections bound the compact chain swept out
by the intervening finite orbit segments, which proves the homology
assertion without a filling behind either boundary.

For clarity, the identification of boundary labels does not assert
that the two hypersurfaces have the same transverse defining
function.  In the attachment coordinates of
Proposition~\ref{st:attachment}, $B_0=\{U=V=0\}$ and the right wedge
has $U,V>0$.  A spacelike hypersurface through $B_0$ is transverse to
both null faces.  If $s\geq0$ is a smooth defining function on that
hypersurface, then
\[
 U=s\,a(s,y),\qquad V=s\,b(s,y),\qquad a(0,y),b(0,y)>0.
\]
Thus $\sqrt{UV}=s\sqrt{ab}$ is a smooth defining function, and the
restriction of $y$ to $B_0$ is its original boundary label.  This also
proves that the graph approaches precisely that boundary and no
other limiting orbit.  The end assertion follows from the rest-chart
graph estimates in Proposition~\ref{st:maximal-slice}.
\end{proof}

Use primes for the geometry of $\Sigma'$:
\[
 \xi=u'n'+X',\qquad \operatorname{tr}_{g'}K'=0,\qquad
 \operatorname{sym}\nabla X'=-u'K'.
\]
Here $u'>0$ in the interior because $\xi$ is future causal and nonzero,
whereas $u'=X'=0$ on $B_0$.  The last identity follows by restricting
$\mathcal L_\xi\mathbf G=0$ to the hypersurface, with the positive
second-fundamental-form convention of Definition~\ref{def:data}.

If $Q>0$, perform the constant duality rotation
\begin{equation}
\label{ri:duality}
 \widehat{\mathbf F}
   =\frac{Q_E}{Q}\mathbf F+\frac{Q_B}{Q}*_{\mathbf G}\mathbf F,
 \qquad
 *_{\mathbf G}\widehat{\mathbf F}
   =\frac{Q_E}{Q}*_{\mathbf G}\mathbf F
       -\frac{Q_B}{Q}\mathbf F.
\end{equation}
If $Q=0$, take $\widehat{\mathbf F}=\mathbf F$.
Let $\mathcal E',\mathcal B'$ be its induced fields on $\Sigma'$.
Their charges are $Q,0$.  Indeed the original spatial restrictions
of $\mathbf F,*\mathbf F$ are $\alpha_2,-\alpha_1$, respectively;
Maxwell closure and Lemma~\ref{ri:orbit-graph} preserve their
integrals.  The rotation also preserves the electromagnetic stress.

\begin{proposition}[Staticity and vanishing of matter]
\label{ri:staticity}
On $\Sigma'$ one has
\[
 K'=0,\qquad X'=0,\qquad \mathcal B'=0.
\]
The Killing field is strictly timelike throughout the open right
wedge, and its metric and rotated Maxwell form have the global form
\begin{equation}
\label{ri:static-product}
 \mathbf G=-(u')^2dT^2+g',\qquad
 \widehat{\mathbf F}=u'\,dT\wedge(\mathcal E')^\flat,
 \qquad u'=\sqrt W.
\end{equation}
The residual non-electromagnetic stress vanishes.  The original
constraint densities $\mu_m,J_m$ consequently vanish on all of
$\Omega$, including $S$.
\end{proposition}

\begin{proof}
The exact potentials of Lemma~\ref{va:potentials} remain exact
after rotation and pullback to the global graph.  In particular,
\begin{equation}
\label{ri:magnetic-potential}
 b'=\iota_\xi(*_{\mathbf G}\widehat{\mathbf F})|_{T\Sigma'}
    =u'(\mathcal B')^\flat+(X'\times\mathcal E')^\flat=d\psi.
\end{equation}
The primitive extends smoothly to $B_0$.  To see this without an
assumption on its periods at the boundary, take a smooth collar
primitive locally, subtract its constant difference from the
interior primitive, and glue; smoothness of the one-form supplies
the local primitives.  Since $\xi=0$ on $B_0$, its differential
vanishes there.  Connectedness of $B_0$, inherited from $S$, makes
$\psi$ one constant $\psi_0$ on the whole boundary.

At infinity $b'=O(r^{-2})$.  Radial integration gives a finite
limit along each ray, and joining two rays along a sphere gives
oscillation at most $Cr^{-1}$.  Hence there is a single constant
$\psi_\infty$ with $\psi-\psi_\infty=O(r^{-1})$.
Since $\operatorname{div}_{g'}\mathcal B'=0$ and its total flux is
zero, the divergence theorem on the portion of $\Sigma'$ inside a
large coordinate sphere gives
\[
 \begin{split}
 \int_{\Sigma'}b'(\mathcal B')\,dV_{g'}
 &=\lim_{R\to\infty}\int_{S_R}
       (\psi-\psi_\infty)g'(\mathcal B',\nu_R)\,dA_{g'}
       -\psi_0\int_{B_0}g'(\mathcal B',\nu)\,dA_{g'}\\
 &=0.
 \end{split}
\]
The outer error is $O(R^{-1})$ and the inner total flux is zero.
This is precisely where a single connected boundary is needed:
zero total magnetic charge would not cancel independently chosen
constants on several components.  Equation~\eqref{ri:magnetic-potential}
therefore yields
\begin{equation}
\label{ri:magnetic-integral}
 \int_{\Sigma'}X'\cdot(\mathcal E'\times\mathcal B')\,dV_{g'}
       =-\int_{\Sigma'}u'|\mathcal B'|^2\,dV_{g'}.
\end{equation}

We check the normalization and sign of the other integral.
Proposition~\ref{va:stationary-data} gives
\[
 \operatorname{Ein}_{\mathbf G}-\mathsf S
       =h_*\xi^\flat\otimes\xi^\flat,
 \qquad h_*\geq0,\qquad\operatorname{supp}h_*\subset\{W=0\}.
\]
Here $h_*$ is supported away from the inner and outer collars.
For the geometric constraint densities used throughout this paper,
$\mathsf S(n',e_i)=2(\mathcal E'\times\mathcal B')_i$.
Also $\mathbf G(\xi,n')=-u'$ and
$\mathbf G(\xi,e_i)=X'_i$.  The mixed Einstein constraint on the
maximal slice is consequently
\begin{equation}
\label{ri:mixed-constraint}
 \nabla^jK'_{ij}
      =2(\mathcal E'\times\mathcal B')_i-h_*u'X'_i.
\end{equation}
There is no $8\pi$ in this identity.  The $8\pi$ factor belongs to
the ADM momentum flux and to conversion from physical matter
densities, not to these geometric densities.

Multiply Equation~\eqref{ri:mixed-constraint} by $X'^i$ and
integrate.  Integration by parts and the Killing equation give
\[
 \int_{\Sigma'} X'^i\nabla^jK'_{ij}\,dV_{g'}
      =-\int_{\Sigma'}K'_{ij}\nabla^jX'^i\,dV_{g'}
      =\int_{\Sigma'}u'|K'|^2\,dV_{g'}.
\]
The inner boundary term vanishes because $X'=0$ there.  At infinity
$X'=O(r^{-q_2})$ and $K'=O(r^{-1-q_2})$ for some $q_2>1/2$,
so its absolute value is $O(R^{1-2q_2})\to0$.  All volume integrals
converge: the gravitational integrand is $O(r^{-2-2q_2})$, the
Maxwell squared norms are $O(r^{-4})$, and $h_*$ has compact
support.  Substituting Equation~\eqref{ri:magnetic-integral} gives
the exact nonnegative identity
\begin{equation}
\label{ri:staticity-integral}
 \int_{\Sigma'}u'
       \bigl(|K'|^2+2|\mathcal B'|^2+h_*|X'|^2\bigr)\,dV_{g'}=0.
\end{equation}
It follows that $K'=\mathcal B'=0$ in the interior and hence at its
boundary by smoothness.

Now $X'$ is a spatial Killing field and vanishes on the whole
boundary.  Its covariant derivatives in boundary tangent
directions are zero.  Skew symmetry of $\nabla X'$ then forces
all remaining components of $\nabla X'$ to vanish there as well.
The Killing identities form the first-order linear system
$\nabla X'=A$, $\nabla A=\operatorname{Rm}_{g'}*X'$ along any
curve.  Its zero initial data at a boundary point give $X'=0$
along every such curve, and connectedness gives $X'=0$ everywhere.
Thus $W=(u')^2>0$ in the interior.  Every orbit meets this slice,
and $W$ is invariant under the flow, so the same conclusion holds
throughout the right wedge.  The support condition now forces
$h_*=0$.  Flowing the orthogonal slice gives
Equation~\eqref{ri:static-product}; invariance and the vanishing of
the tangential restriction of $\widehat{\mathbf F}$ give its stated
form.  Restriction back to $z=0$ proves $\mu_m=J_m=0$ in the
original open exterior.  Smoothness of the original data extends
these equalities to $S$.
\end{proof}

\subsection{The mass and boundary of the static orbit metric}
\label{ri:mass-area}

On the open orbit space, the now positive function $W$ gives
\begin{equation}
\label{ri:orbit-metric}
 g'=h_s=g+W^{-1}X_g^\flat\otimes X_g^\flat\geq g.
\end{equation}
This is an equality of metrics after orbit projection in the
interior; different boundary defining functions are allowed.
The induced metric of $B_0$ equals the original metric on $TS$ by
the attachment.  Thus $|B_0|_{g'}=A$.
The fixed set $B_0$ is totally geodesic in spacetime: its second
fundamental vector is fixed by the differential of every Killing
flow, while that differential acts as a nontrivial boost in the
normal two-plane and has no fixed vector there.  In particular
$B_0$ is minimal in $g'$.  The constraint and static equations give
\begin{equation}
\label{ri:static-equations}
 R_{g'}=2|\mathcal E'|_{g'}^2,\qquad
 \operatorname{div}_{g'}\mathcal E'=0,\qquad
 \Delta_{g'}u'=|\mathcal E'|_{g'}^2u'.
\end{equation}
For the last equation, the Ricci component in the future unit
normal direction of a static metric is $(u')^{-1}\Delta_{g'}u'$.
Electrovacuum has trace-free stress and that component of its
Ricci tensor is $|\mathcal E'|^2$.

\begin{lemma}[The original invariant mass is the static ADM energy]
\label{ri:static-mass}
The metric $g'$ is complete up to $B_0$, has one strongly
asymptotically flat end, and in its stationary rest coordinates
\[
 g'_{ij}-\delta_{ij}=O_2(r^{-1}),\qquad
 \mathcal E'=O_1(r^{-2}),\qquad E_{\mathrm{ADM}}(g')=m.
\]
Its electric charge is $Q$ and its magnetic charge is zero.
\end{lemma}

\begin{proof}
Smoothness through $B_0$, compactness of the core and completeness
of the end follow from Proposition~\ref{st:maximal-slice} and
Lemma~\ref{ri:orbit-graph}.  The symbol bounds for the stationary
end in Proposition~\ref{st:stationary-end}, together with
Equation~\eqref{ri:orbit-metric}, give the stated strong
asymptotic flatness and field bounds.  Fluxes were checked in the
proof of Proposition~\ref{ri:staticity}.

We spell out why the mass is that of the original data.  In the
unimproved stationary coordinates the spacetime perturbation of
its limiting constant Minkowski metric has differentiated decay
$O_2(r^{-q_0})$, $q_0>1/2$.  It is electrovacuum near infinity.
The skew-slot flux in Lemma~\ref{en:adm-flux} therefore has
spacetime divergence
\[
 O(r^{-2-2q_0})+O(r^{-4}).
\]
Stokes' Theorem between cuts in the original asymptotic plane and
cuts in a plane orthogonal to the limiting Killing direction has
error
$O(R^{1-2q_0})+O(R^{-1})\to0$: the connecting chains have volume
$O(R^3)$ and remain at radii comparable to $R$.  The same estimate
permits replacing comparable large cuts by coordinate spheres.
With the positive-$K$ convention, Lemma~\ref{en:adm-flux} pairs
a translation $(b^0,b^i)$ with $b^0E+\sum_i b^iP_i$.  Here the
limiting unit Killing translation furnished by
Proposition~\ref{va:stationary-data} is
\[
 b^0=\frac E m,\qquad b^i=-\frac{P_i}{m}.
\]
Its flux on the original section is therefore
\[
 b^0E+\sum_i b^iP_i
       =\frac{E^2-|P|^2}{m}=m.
\]
On the orthogonal rest plane that same flux is its ADM energy.
Its spatial metric and $h_s$ differ by the quadratic shift term,
which is $O_1(r^{-2q_0})$ in this chart and contributes
$O(R^{1-2q_0})$ to the ADM flux.

Finally the stationary harmonic correction has spatial
coordinate displacement $O_2(r^{1-q_1})$, $q_1>1/2$, as in
Proposition~\ref{st:stationary-end}.  Its leading metric change
is a symmetric flat gradient $2\partial_{(i}a_{j)}$.
The ADM divergence of this term is
$\Delta a_i-\partial_i\operatorname{div}a$, whose divergence
vanishes identically.  Extend $a$ smoothly across the inside of a
coordinate sphere; the flux of this leading term is exactly zero
by the Euclidean divergence theorem.  Products of coordinate and
metric errors have flux $O(R^{1-2q_1})\to0$, with a smaller decay
exponent chosen if necessary.  The energy is consequently still
$m$ in the strong asymptotic coordinates.  This argument uses the
original invariant ADM flux, not the mass of a deformation.
\end{proof}

\begin{proposition}[Identification of the static exterior]
\label{ri:rn-identification}
There is an orientation-preserving isometry from $(\Sigma',g')$
onto the canonical exterior Reissner--Nordstr\"om slice of mass $m$
and charge magnitude $Q$, including its boundary.  Under this
isometry, with $r$ the areal radius,
\begin{equation}
\label{ri:rn-static-data}
 g'=f^{-1}dr^2+r^2\sigma_2,\qquad
 \mathcal E'=\frac Q{r^2}\sqrt f\,\partial_r,\qquad
 u'=\sqrt f,\qquad
 f=1-\frac{2m}{r}+\frac{Q^2}{r^2},\quad r\geq r_h.
\end{equation}
\end{proposition}

\begin{proof}
Put $N=u'$.  We first verify the electrostatic system, including its
potential and asymptotic normalization.  By Lemma~\ref{va:potentials}
and the constant duality rotation, the one-form
$\iota_\xi\widehat{\mathbf F}=N(\mathcal E')^\flat$ has a global
primitive.  Choose its negative, denoted by $\Psi$, and add a
constant so that $\Psi\to0$ at infinity.  This normalization is
possible because $d\Psi=O(r^{-2})$: radial integration gives a
limit along each ray, and the oscillation on a sphere is $O(r^{-1})$.
Thus
\[
 d\Psi=-N(\mathcal E')^\flat,\qquad
 \widehat{\mathbf F}=d\Psi\wedge dT,\qquad \Psi=O(r^{-1}).
\]
The primitive extends smoothly to $B_0$, as in the proof of
Proposition~\ref{ri:staticity}.  Its differential vanishes there
because $\xi=0$, and connectedness makes its boundary value a
single constant $\Psi_0$.  This is a scalar Killing potential;
no global vector potential or simple connectivity is needed.

The spatial Ricci tensor of $-N^2dT^2+g'$ is
$\operatorname{Ric}_{g'}-N^{-1}\nabla^2N$.  The spatial component
of the purely electric Maxwell Ricci tensor is
$|\mathcal E'|^2g'-2(\mathcal E')^\flat\otimes(\mathcal E')^\flat$.
Together with Equation~\eqref{ri:static-equations} and
$\operatorname{div}_{g'}\mathcal E'=0$, this gives
\[
 \begin{aligned}
 N\operatorname{Ric}_{g'}
   &=\nabla^2N-\frac2N\,d\Psi\otimes d\Psi
                        +\frac1N|d\Psi|^2g',\\
 \Delta_{g'}N&=\frac1N|d\Psi|^2,\qquad
 \Delta_{g'}\Psi=\frac1N\langle dN,d\Psi\rangle.
 \end{aligned}
\]
These signs agree with the convention
$\widehat{\mathbf F}=d\Psi\wedge dT$.

We next establish the lapse expansion required by electrostatic
uniqueness.  Proposition~\ref{st:stationary-end} gives
$N=1+O_j(r^{-1})$, $g'-\delta=O_j(r^{-1})$, and
$\mathcal E'=O_j(r^{-2})$ for every fixed $j$, with scaled
H\"older bounds.  Hence the lapse equation implies
$\Delta_\delta(N-1)=O_j(r^{-4})$.
Multiply $N-1$ by an end cutoff and extend it by zero to
$\mathbb R^3$.  Its Euclidean Laplacian $b$ is integrable and is
$O_j(r^{-4})$ off a compact set.  The Newton potential
$-(4\pi)^{-1}\int b(y)|x-y|^{-1}\,dy$ equals this cutoff
function, since their difference is an entire harmonic function
tending to zero.  Subtracting the monopole term and splitting the
integral into $|y|<r/2$, $|y|\asymp r$, and $|y|>2r$ gives
an $O(r^{-2}\log r)$ remainder.  Scaled Poisson estimates give
the corresponding first two derivative bounds.  Consequently,
for a constant $\mu$,
\[
 N=1-\frac\mu r+O_2(r^{-2}\log r)
       =1-\frac\mu r+o_2(r^{-1}).
\]

For completeness, this coefficient is the original invariant mass.
Let $\mathcal G_{ij}=\operatorname{Ric}_{g',ij}
-\tfrac12R_{g'}g'_{ij}$.  Expanding about the Euclidean metric,
the linearized contracted Bianchi identity gives
\[
 \partial_j\left(\mathcal G^{\rm lin}_{ij}x^i
       +\tfrac12(\partial_i g'_{ij}-\partial_j g'_{ii})\right)=0.
\]
Extend the metric perturbation smoothly through a fixed ball.
The flux of this divergence-free vector field is zero.  The
nonlinear terms have size $O(r^{-4})$
and give an $O(r^{-1})$ error in the resulting flux identity.
Thus, with $n_\delta$ and $dA_\delta$ Euclidean,
\[
 E_{\mathrm{ADM}}(g')
 =-\frac1{8\pi}\lim_{R\to\infty}
       \int_{S_R}\mathcal G_{ij}x^i n_\delta^j\,dA_\delta.
\]
The electrostatic equations give
$\mathcal G=N^{-1}\nabla^2N
-2(\mathcal E')^\flat\otimes(\mathcal E')^\flat
=\partial^2N+O(r^{-4})$ in these coordinates.
The last flux is therefore $-8\pi\mu+o(1)$, proving
$\mu=E_{\mathrm{ADM}}(g')=m$ by Lemma~\ref{ri:static-mass}.
Also, with $\nu$ pointing from $B_0$ toward the end, integration
of the lapse equation gives
\[
 4\pi\mu=\int_{B_0}\partial_\nu N\,dA
              +\int_{\Sigma'}N|\mathcal E'|^2\,dV>0.
\]
Here $\partial_\nu N=\kappa>0$ follows from the bifurcation
attachment and the normalization of $\xi$.

We can now apply the connected-horizon electrostatic uniqueness
theorem of Borghini, Cederbaum, and Cogo
\cite[Theorem~3.1 and Proposition~6.1]{BorghiniCederbaumCogo2025}.
The spatial manifold is connected, oriented, one-ended and complete
up to its compact connected boundary by Lemma~\ref{ri:static-mass}.
The metric, lapse and potential are smooth up to that actual
boundary, $N>0$ inside, $N=0$ and $|dN|=\kappa>0$ on $B_0$,
and $\Psi$ is constant there.  The displayed electrostatic
system and asymptotics are exactly the required ones.  The
additional charge expansion imposed in that theorem when
$\mu=0$ is irrelevant because $\mu=m>0$.
It follows that the entire system is the subextremal
Reissner--Nordstr\"om system with mass $m$ and some signed
charge $q$, including its normalized lapse and potential
$N=\sqrt{1-2m/r+q^2/r^2}$ and $\Psi=q/r$.
Its isometry extends smoothly to the nondegenerate boundary.

The sign is fixed by the actual flux, not by a choice of potential
convention: $\mathcal E'=-N^{-1}\nabla\Psi$ for our convention,
so the charge of this model is $q$.  Lemma~\ref{ri:static-mass}
therefore gives $q=Q$.  This also covers $Q=0$.  More directly,
in that case integration of
$\operatorname{div}(\Psi N^{-1}\nabla\Psi)
=N^{-1}|d\Psi|^2$ gives
\[
 \int_{\Sigma'}N^{-1}|d\Psi|^2\,dV=4\pi\Psi_0 Q=0.
\]
The boundary expression is regular because
$N^{-1}d\Psi=-(\mathcal E')^\flat$; the term at infinity tends
to zero.  Thus $\Psi=0$, and the zero-charge case of the same
uniqueness theorem is Schwarzschild.

The identified horizon radius is
$m+\sqrt{m^2-Q^2}=r_h$.  The boundary isometry is smooth in
proper-distance collars, where $N$ vanishes simply.  If the
isometry initially reverses orientation, compose it with a
reflection of the round sphere.  That reflection fixes the radial
lapse, potential and electric field, and makes the resulting
identification orientation preserving.  This proves all three
identifications in Equation~\eqref{ri:rn-static-data} on the
closed exterior.
\end{proof}

\subsection{Both Maxwell charges and the original horizon section}
\label{ri:original-subsection}

With the spacetime orientation fixed in Section~\ref{sec:statements},
an orthonormal coframe
$e^0,e^1,e^2,e^3$ with volume $e^0\wedge e^1\wedge e^2\wedge e^3$
satisfies
\[
 *(e^0\wedge e^1)=-e^2\wedge e^3,\qquad
 *(e^2\wedge e^3)=e^0\wedge e^1.
\]
Equations~\eqref{ri:static-product} and
\eqref{ri:rn-static-data} therefore give
\[
 \widehat{\mathbf F}=\frac Q{r^2}dT\wedge dr,
 \qquad *\widehat{\mathbf F}=-Q\,dA_{\sigma_2}.
\]
For $Q>0$ the inverse of Equation~\eqref{ri:duality} is
\begin{equation}
\label{ri:inverse-duality}
 \mathbf F=\frac{Q_E}{Q}\widehat{\mathbf F}
              -\frac{Q_B}{Q}*\widehat{\mathbf F}
       =\frac{Q_E}{r^2}dT\wedge dr+Q_B\,dA_{\sigma_2}.
\end{equation}
For $Q=0$ all these forms vanish and the same conclusion holds.
In particular,
\[
 *\mathbf F=\frac{Q_B}{r^2}dT\wedge dr-Q_E\,dA_{\sigma_2}.
\]
Contraction with the future static unit normal $f^{-1/2}\partial_T$
gives the two outward fields $Q_Er^{-2}\sqrt f\,\partial_r$ and
$Q_Br^{-2}\sqrt f\,\partial_r$.  Thus neither charge sign has been
changed by the identification.

\begin{proposition}[Global realization of the original data]
\label{ri:original-embedding}
The original stationary section $z=0$ gives a smooth
orientation-preserving spacelike embedding of all of $\Omega$
into one Reissner--Nordstr\"om exterior together with its future
horizon or bifurcation sphere.  It induces the original
$(g,K,\mathcal E,\mathcal B)$, including their values on $S$,
and its end approaches the corresponding spatial infinity.
\end{proposition}

\begin{proof}
The global product in Equation~\eqref{ri:static-product} and
Proposition~\ref{ri:rn-identification} identify the right wedge
with the static Reissner--Nordstr\"om exterior, with the form in
Equation~\eqref{ri:inverse-duality}.  The original section is a
global graph there by Lemma~\ref{ri:orbit-graph}.  All original
tensors are already induced by $\mathbf G,\mathbf F$ on that
section; in particular its second fundamental form is the
original $K$, not the zero second fundamental form of $\Sigma'$.
Its unit normal need not be the static normal.  The contraction
identities defining the original fields consequently recover
both full fields even when they are nonradial or nonparallel on
that section.

It remains to verify smoothness and injectivity at the original
boundary.  On the open original orbit base write
\begin{equation}
\label{ri:time-connection}
 \mathbf G=-W(dz-A_s)^2+h_s,\qquad
 A_s=W^{-1}X_g^\flat,\qquad dT=dz-A_s.
\end{equation}
The last equality is global: $T$ is the Killing-flow coordinate
based on the orthogonal global section $\Sigma'$, so $T-z$ is a
single-valued function of the orbit label.  On the original
section it says $dT=-A_s$.
Let $r_*$ be the Reissner--Nordstr\"om tortoise coordinate,
$dr_*/dr=f^{-1}$.  Since $f'(r_h)=2\kappa$ and $W=f(r)$,
\begin{equation}
\label{ri:tortoise}
 r_*=(2\kappa)^{-1}\log W+a(W)
\end{equation}
for a smooth function $a$ near $W=0$.  To verify this expansion,
use the smooth inverse $r=r(W)$ of $f$ near $r_h$ and subtract
$(2\kappa W)^{-1}$ from $dr_*/dW$; the difference is smooth.
The normalization of $\kappa$ is the one from
Proposition~\ref{va:boundary}, equivalently the surface gravity
of the now normalized static Killing field.

If $u|_S=0$, Proposition~\ref{st:attachment} gives smooth
$A_s$ and a smooth defining function $\sqrt W$ in the original
collar.  Hence $T|_{z=0}$ extends smoothly there.  The orbit
metric is a smooth nondegenerate metric in that collar, and the
base isometry extends smoothly by its normal geodesic collars;
its angular boundary label is a diffeomorphism.  The exterior
Kruskal coordinates
\[
 \mathsf U=-e^{\kappa(r_*-T)},\qquad
 \mathsf V=e^{\kappa(r_*+T)}
\]
are, by Equation~\eqref{ri:tortoise}, smooth multiples of
$\sqrt W$ on this section.  Both vanish simply in its inward
normal parameter, while the angular variables parametrize $S$.
Their derivative vector is a nonzero spacelike direction in
the normal two-plane, since the two nonzero factors have
opposite signs.  This gives a smooth immersion through the
bifurcation sphere.

The exterior in these Kruskal coordinates has $\mathsf U<0<\mathsf V$;
the attachment coordinates of Section~\ref{st:section} have $U,V>0$.

If $u|_S>0$, the other part of Proposition~\ref{st:attachment}
gives $W$ as a smooth original boundary defining function and
\[
 X_g^\flat=(2\kappa)^{-1}dW+W\beta,
 \qquad A_s=(2\kappa)^{-1}d\log W+\beta,
\]
where $\beta$ is smooth in $(W,y)$.  Equations~\eqref{ri:time-connection}
and \eqref{ri:tortoise} show that the advanced time
\[
 v=T+r_*
\]
has smooth differential $dv=-\beta+a'(W)dW$ on the original
section, and hence has a finite smooth extension to $S$.
In these coordinates the spacetime metric is
\[
 -f\,dv^2+2\,dv\,dr+r^2\sigma_2.
\]
The sign of the logarithmic cancellation therefore selects the
future horizon, where $r=r_h$ and $v$ is finite.

Here the angular map also is smooth in the original defining
function $W$, a fact that requires more than smoothness in
$\sqrt W$.  Put $s=\sqrt W$.  The formula for $X_g^\flat$ shows
that
\[
 h_s=g+W^{-1}\bigl((2\kappa)^{-1}dW+W\beta\bigr)^2
\]
extends as a smooth nondegenerate tensor in $(s,y)$ invariant
under $s\mapsto-s$.  At $s=0$ its normal coefficient is
$\kappa^{-2}$ and its tangential metric is $g|_{TS}$.
Reflection is consequently a local isometry fixing the
boundary.  In normal geodesic collars the angular labels are
constant along the normal geodesics and invariant under this
reflection.  The angular components of the base isometry are
thus smooth even functions of $s$.  Taylor's formula with
parameters makes every such function smooth in $s^2=W$ on the
one-sided collar.  The radial coordinate $r=r(W)$ is smooth as
well.  We have obtained a smooth map in the regular coordinates
$(v,r,\text{angles})$, and $dr/dW=(2\kappa)^{-1}$ at $S$ is
nonzero.  Together with the boundary angular diffeomorphism
this proves immersion at the future horizon and shows that
$S$ maps onto a full horizon cross-section.

In either case the limiting pullback of the ambient metric is
the original smooth positive definite $g$, so the immersion is
spacelike up to $S$.  The induced future normal is a smooth
algebraic function of its first derivatives and the ambient
metric; it therefore extends smoothly as well.  The second
fundamental form and the two field contractions agree in the
open exterior and extend by smoothness, proving their equality
with the original tensors on $S$.

The open graph is injective, its boundary angular map is
injective, and the boundary image is disjoint from the open
exterior.  Finally its end escapes properly: the base radius
tends to infinity and the rest-chart graph has the original
strictly spacelike asymptotic plane slope, up to sublinear
errors.  Its image approaches spatial infinity in that end.
Compactness of the remaining core and the boundary collar
then make the injective immersion proper, hence an embedding.
The spatial orientation was fixed in
Proposition~\ref{ri:rn-identification}; flowing the future
normal preserves the corresponding spacetime orientation.
This proves the stated orientation and all the required
original-data matching.
\end{proof}

\begin{proof}[Proof of Theorem~\ref{thm:rigidity}]
The equality variation and stationary construction in
Propositions~\ref{va:stationary-data}, \ref{va:boundary}, and
\ref{st:maximal-slice}, together with Lemma~\ref{va:potentials} apply to the
original equality data.  Proposition~\ref{ri:staticity} proves
static electrovacuum in their entire right wedge.
Lemma~\ref{ri:static-mass} preserves the original invariant mass,
and Proposition~\ref{ri:rn-identification} identifies its orbit
metric, lapse, and rotated electric field.  The inverse duality
formula~\eqref{ri:inverse-duality} recovers the prescribed two
charge signs.  Proposition~\ref{ri:original-embedding} then
realizes the entire original data, with smooth future normal
and fields at its horizon boundary.  The vanishing of original
non-electromagnetic matter was proved in
Proposition~\ref{ri:staticity}.
\end{proof}

\subsection{Converse and three families of sharp examples}
\label{ri:examples-subsection}

\begin{proof}[Proof of Theorem~\ref{thm:converse}]
Write $Q_*^2=q_E^2+q_B^2$ and
$R_+=M+\sqrt{M^2-Q_*^2}$.  The degenerate induced metric on the
future horizon is $R_+^2\sigma_2$, with its null direction along
the generators.  Every full smooth cross-section, including the
bifurcation sphere, therefore has area $4\pi R_+^2$.  The assumed
original cut condition implies
\[
 A_{\min}(S)=|S|_g=4\pi R_+^2.
\]
The converse assumes that the actual invariant ADM mass is $M$
and that the actual outward fluxes are $q_E,q_B$.  Substitution
now gives the required equality.  The examples below verify
those assumptions for three explicit types of slicing.
\end{proof}

\begin{proposition}[Static, non-time-symmetric, and boosted examples]
\label{ri:sharp-examples}
For every $M>\sqrt{q_E^2+q_B^2}$ there are admissible equality
data of each of the following kinds: static data; data with
$K\not\equiv0$; and data with nonzero ADM momentum.  In the
last family, for every sufficiently small nonzero vector $v$,
the actual ADM energy and momentum are
\[
 E=\gamma M,\qquad P=\gamma Mv,\qquad
 \gamma=(1-|v|^2)^{-1/2},
\]
with the positive-$K$ convention of this paper.  In all three
families the actual outward charges are $q_E,q_B$ and the
original minimum enclosing area is $4\pi R_+^2$.
\end{proposition}

\begin{proof}
Use the converse parameters $M,q_E,q_B$ and set $Q_*^2=q_E^2+q_B^2$. Let
\[
 f(r)=1-\frac{2M}{r}+\frac{Q_*^2}{r^2},\qquad
 g_0=f^{-1}dr^2+r^2\sigma_2,
 \qquad r\geq R_+.
\]
Start with the static exterior down to the bifurcation sphere.
It is smooth there in proper distance or Kruskal coordinates;
its future normal and both induced Maxwell fields are smooth,
$K=0$, and its boundary is a connected minimal sphere with
$\theta_+=0$.  Its metric is complete with that boundary included.
Consider also time graphs $T=h(y)$ in areal Cartesian variables
$y=r\omega$ of the following two types:
\begin{enumerate}
 \item $h=\varepsilon h_0(r)$, where $h_0$ is smooth with compact
 support in $(R_+,\infty)$, and has a critical point $r_0$ with
 $h_0''(r_0)\ne0$;
 \item $h=\chi(r/L)\,v\cdot y$, where $\chi=0$ for $r/L\leq1$,
 $\chi=1$ for $r/L\geq2$, $L$ is sufficiently large, and $|v|$
 is sufficiently small and nonzero.
\end{enumerate}
All graphs agree with the static slice near $S$.  Their induced
metric is
\begin{equation}
\label{ri:graph-metric}
 g_h=g_0-f\,dh\otimes dh.
\end{equation}
For the radial graph it is
$(f^{-1}-f(h')^2)dr^2+r^2\sigma_2$ and is positive for small
$\varepsilon$.  For the cutoff graph, bounded derivatives of
$\chi$ give $f|dh|_{g_0}^2\leq C|v|^2$.  Choose $|v|$ such that
this is at most a fixed $\eta<1$.  Then
\begin{equation}
\label{ri:metric-domination}
 g_h\geq(1-\eta)g_0.
\end{equation}
These estimates prove spacelikeness and completeness, and the
unchanged inner collar supplies smoothness of the graph and its
future normal at the bifurcation sphere.

These are electrovacuum initial data with exactly the
normalizations of Definition~\ref{def:data}.  One can check
this directly in the static coordinates: the mixed diagonal
components of the Einstein tensor are
\[
 \frac{rf'+f-1}{r^2},\quad
 \frac{rf'+f-1}{r^2},\quad
 \frac{f''}{2}+\frac{f'}r,\quad
 \frac{f''}{2}+\frac{f'}r,
\]
which equal $r^{-4}(-Q_*^2,-Q_*^2,Q_*^2,Q_*^2)$, the components
of
$2(F_{ac}F_b{}^c-\frac14\mathbf G_{ab}F_{cd}F^{cd})$.
The forms
\[
 F=\frac{q_E}{r^2}dT\wedge dr+q_BdA_{\sigma_2},\qquad
 *F=\frac{q_B}{r^2}dT\wedge dr-q_EdA_{\sigma_2}
\]
are closed.  Gauss--Codazzi with
$K(U,V)=\mathbf G(\nabla_U n,V)$ gives
$\mu=|\mathcal E|^2+|\mathcal B|^2$ and
$J=2(\mathcal E\times\mathcal B)^\flat$ on every graph; the
spatial pullbacks of the closed forms give both divergence
constraints.  Thus $\mu_m=J_m=0$ and the matter DEC holds.
For every graph sphere $r=R$, $dr$ vanishes on its tangent
space, so directly
\begin{equation}
\label{ri:example-flux}
 \frac1{4\pi}\int_{r=R}F=q_B,\qquad
 -\frac1{4\pi}\int_{r=R}*F=q_E.
\end{equation}
The orientation is toward increasing $r$.  These are precisely
the induced magnetic and electric fluxes, respectively.
Closure identifies them with the fluxes on the actual end
coordinate spheres, even when those coordinate spheres are
not the surfaces $r=R$.

We next check the full cut condition in the original graph
metric.  On the product $[R_+,\infty)\times\mathbb S^2$ let
\[
 \alpha=R_+^2dA_{\sigma_2}.
\]
It is closed and has comass $R_+^2/r^2\leq1$ in the static and
radial graph metrics, since their radial and angular blocks are
orthogonal and the angular metric is $r^2\sigma_2$.
For the cutoff graph, Equation~\eqref{ri:metric-domination}
shows that its comass is at most
$R_+^2/((1-\eta)r^2)$ in the changed region.  Choose
$L>R_+/\sqrt{1-\eta}$.  On the unchanged inner region the
previous exact bound applies, so the comass is at most one
everywhere on this graph too.

If $\Gamma=\partial D$ is any full cut as in
Definition~\ref{def:full-cut}, truncate $D$ by a large sphere
outside $\Gamma$.  Stokes' Theorem for this compact region,
with all its inner boundary components and coincident obstacle
portions counted, gives
\[
 \int_\Gamma\alpha=\int_{r=R}\alpha=4\pi R_+^2.
\]
The comass bound yields
$|\Gamma|_{g_h}\geq4\pi R_+^2=|S|_{g_h}$.
It covers disconnected cuts and cuts meeting the obstacle.
Taking $D=\Omega$ shows the infimum equals this value.  No
extension or filling behind the horizon has entered the
calculation.

The boundary expansion remains zero on the unchanged collar.
Every coordinate sphere with $r>R_+$ has strictly positive
$\theta_+$ on these examples.  On the static slice its value
is $2\sqrt f/r$.  For the compact radial change, the value
varies continuously with the height in $C^2$ on its fixed
compact support, so sufficiently small $\varepsilon$ preserves
strict positivity; outside that support it has the static
value.  For the cutoff graph choose $L$ large enough that
$2\sqrt f\geq1$ for $r\geq L$.  On any annulus with radii
comparable to $R\geq L$, rescale spatial coordinates and time
by $R$.  The background metric and its derivatives are
uniformly bounded and converge to the Minkowski ones as
$R\to\infty$, and the rescaled height has $C^2$ norm at most
$C|v|$, with a constant independent of $R$ and of whether the
annulus meets the cutoff.  The induced normal, mean curvature,
and tangential $K$ trace depend smoothly on these derivatives
as long as the graph stays uniformly spacelike.  Consequently
\[
 \left|r\theta_+(r)-2\sqrt{f(r)}\right|\leq C'|v|
 \qquad(r\geq L)
\]
pointwise on each sphere.  Taking $C'|v|<1/2$ proves the claim.

To deduce the required outermostness, suppose a compact smooth
enclosing surface entirely in the open graph had
$\theta_+\leq0$.  At a point of largest $r$ it is tangent to
the coordinate sphere, with the same outward normal: the
radially outward side is in the component containing infinity.
The inner surface's local graph lies on the inner side of the
coordinate sphere.  In coordinates on their common tangent
plane, the second derivative test and the mean-curvature
principal part $-\Delta$ give
$H_{\mathrm{surface}}\geq H_{\mathrm{sphere}}$ at contact.
Their tangential $K$ traces are equal there because their
tangent planes coincide.  Hence its expansion is at least
the strictly positive sphere expansion, a contradiction.
The argument applies to a largest-radius point among all
components, so it also excludes disconnected enclosing weakly
outer-trapped surfaces.

The radial example has nonzero original second fundamental
form.  At its critical radius the future normal is the static
normal, and the graph Hessian formula gives
\[
 K_{rr}(r_0)=\sqrt{f(r_0)}\,h''(r_0)\ne0.
\]
Thus these examples are not time-symmetric.

For completeness we calculate the ADM fluxes, including the
momentum sign of the boosted examples, directly.  The static
metric in areal Cartesian coordinates is
\[
 (g_0)_{ij}=\delta_{ij}+\frac{2M}{r}\frac{y_i y_j}{r^2}
                   +O_2(r^{-2}).
\]
Its energy flux is $M$ and its momentum flux is zero, as also
follows from the calculation below at $v=0$.  The compact
radial graph has exactly the same end data.

Rotate the spatial axes for the cutoff graph so $v=\beta e_3$,
where $\beta$ may have either sign.  On its uncut tail set
\[
 T=\beta\gamma x^3,\qquad
 y=(x^1,x^2,\gamma x^3),\qquad
 \gamma=(1-\beta^2)^{-1/2},\qquad d=\gamma^2-1.
\]
The induced Minkowski metric in $x$ is Euclidean.  Let
$A_0=\operatorname{diag}(1,1,\gamma^2)$,
$w=A_0x$, and $r^2=x\cdot A_0x$.  The leading induced metric is
\begin{equation}
\label{ri:boost-metric}
 (g_h)_{ij}=\delta_{ij}+H_{ij}+O_2(|x|^{-2}),\qquad
 H_{ij}=2M\left(\frac{d\,\delta_{i3}\delta_{j3}}r
                     +\frac{w_iw_j}{r^3}\right).
\end{equation}
Because $\partial_i r=w_i/r$, differentiation gives
\[
 \partial_jH_{ij}-\partial_iH_{jj}
 =\frac{2M}{r^3}
  \bigl[(\operatorname{tr}A_0+d)w_i-(A_0w)_i
                     -d\delta_{i3}w_3\bigr]
 =\frac{4M\gamma^2x_i}{r^3}.
\]
On $|x|=R$ put $\mu=x^3/R$ and
$\rho^2=1+d\mu^2$, so $r=R\rho$.  The $16\pi$ flux therefore is
\begin{equation}
\label{ri:boost-energy}
 E=\frac{M\gamma^2}{2}
        \int_{-1}^1(1+d\mu^2)^{-3/2}\,d\mu
   =\gamma M,
\end{equation}
where the primitive is $\mu/\sqrt{1+d\mu^2}$.

To compute $K$ with its sign, the leading spacetime perturbation
in $(T,y)$ is $h_{00}=2M/r$, $h_{0a}=0$,
$h_{ab}=2My_ay_b/r^3$.  Its required connection coefficients are
\[
 \Gamma^0_{0a}=\frac{My_a}{r^3},\qquad
 \Gamma^a_{00}=\frac{My_a}{r^3},\qquad
 \Gamma^a_{bc}
   =M\left(\frac{2\delta_{bc}y_a}{r^3}
                     -\frac{3y_ay_by_c}{r^5}\right),
\]
up to $O_1(r^{-3})$; the remaining mixed time coefficients at
this order are zero.  Put $F_0=T-\beta y^3$ and
$\sigma=\sqrt{-\mathbf G^{-1}(dF_0,dF_0)}
=\gamma^{-1}+O(r^{-1})$.  The future normal is
$n=-\sigma^{-1}\nabla F_0$.  For the constant tangent vectors
$e_i$ to the plane this gives
\[
 K_{ij}=-\sigma^{-1}\operatorname{Hess}F_0(e_i,e_j)
  =\gamma(\Gamma^0_{\alpha\lambda}-\beta\Gamma^3_{\alpha\lambda})e_i^\alpha e_j^\lambda
       +O_1(R^{-3}).
\]
Writing $z=x^3$, substitution yields
\begin{equation}
\label{ri:boost-second-form}
 K_{ij}=M\beta\gamma^2
 \left[
  \frac{x_i\delta_{j3}+x_j\delta_{i3}
                    -2z\delta_{ij}-dz\delta_{i3}\delta_{j3}}{r^3}
       +\frac{3zw_iw_j}{r^5}
 \right]+O_1(R^{-3}).
\end{equation}
The trace at this order is
\[
 \tau=M\beta\gamma^2
       \left[-\frac{(4+d)z}{r^3}+\frac{3z|w|^2}{r^5}\right]
                +O(R^{-3}).
\]
Contracting with the Euclidean outward normal $x/R$ gives
\[
 R^2(K_{3j}-\tau g_{3j})\frac{x^j}{R}
  =M\beta\gamma^2
     \frac{1+(3+4d)\mu^2}{(1+d\mu^2)^{5/2}}+O(R^{-1}).
\]
Consequently the actual $8\pi$ momentum flux is
\begin{equation}
\label{ri:boost-momentum}
 P_3=\frac{M\beta\gamma^2}{4}
      \int_{-1}^1\frac{1+(3+4d)\mu^2}{(1+d\mu^2)^{5/2}}\,d\mu
      =\gamma M\beta.
\end{equation}
Indeed an antiderivative is
$\mu[1+(1+2d)\mu^2]/(1+d\mu^2)^{3/2}$, whose endpoint
difference is $4/\gamma$.  The two transverse momentum
components vanish by reflection in the corresponding
coordinate.  Terms quadratic in asymptotic errors and the
charge-squared terms have integrated error $O(R^{-1})$ and do
not alter either flux.  Rotating back gives $P=\gamma Mv$.
Thus the actual invariant mass is
$\sqrt{E^2-|P|^2}=M$, and the momentum is nonzero when $v\ne0$.

Equations~\eqref{ri:boost-metric} and
\eqref{ri:boost-second-form} also verify
$g_h-\delta=O_2(R^{-1})$ and $K=O_1(R^{-2})$ in the actual
Euclidean end chart.  Pulling back the smooth Maxwell forms
and contracting with the uniformly timelike graph normal
gives both fields $O_1(R^{-2})$.  Their stress, and hence
$\mu$ and $|J|$, are $O(R^{-4})$ and integrable.  The compact
core causes no integrability issue.  Together with
Equation~\eqref{ri:example-flux}, the cut calibration, and the
expansion barrier, these checks establish every initial-data
and connected rigidity-subclass hypothesis for all three
families.  The converse just proved supplies equality.
\end{proof}
